\subsection{A cubic Kummer obstruction at the prime 109}

\begin{statement}{E.28}{Proved}{A nonempty cubic abelian fiber with no class-two lifts}
In the pro-$3$ version of E.24, consider the toric section class
\[
 (\kappa(6),\kappa(109))
 \in H^1(\mathbb Q,\mathbb Z_3(1)^2).
 \tag{E28.1}
\]
Its fiber of geometrically abelian sections is nonempty and is
a torsor under
$\mu_3^3\times\mathcal R_{\mathcal I}(\mathbb Z_3)\times\mathbb Z_3$.
No point of this fiber lifts to a section of the curve's own
class-two pro-$3$ quotient. No such lift exists for the intrinsic
metabelian quotient, the intrinsic marked-depth quotient of E.25,
or the inherited class-two pro-$3$ quotient either.
\end{statement}
\begin{proof}
Let $F_3$ be the free pro-$3$ group on the base loops $x,y$,
and put $\mathcal H=F_3/\gamma_3F_3$.
In its class-two Lie coordinates over $\mathbb Z_3$ write
$X=\log x$, $Y=\log y$, $C=[X,Y]$.
The Baker--Campbell--Hausdorff formula has only the denominator
two, which is a unit in $\mathbb Z_3$.
The image $H$ of $N_3$ consists exactly of
\[
 \exp(3aX+3bY+cC),\qquad a,b,c\in\mathbb Z_3.
 \tag{E28.2}
\]
The exponent-sum condition proves containment, and the elements
$x^3,y^3,[x,y]\in N_3$ generate the full lattice in (E28.2).
In coordinates $(a,b,c)$ its multiplication is
\[
 (a,b,c)(a',b',c')=
 \left(a+a',b+b',
 c+c'+\frac92(ab'-ba')\right).
 \tag{E28.3}
\]
In particular the commutator is $(0,0,9(ab'-ba'))$.
The coordinates $a,b$ are the two coordinate-unit homology
coordinates of $U_3$; the covering pushforward to base
homology is $(3a,3b)$.

The normalized Galois action sends $x$ to $x^\chi$ and
$y$ to $f_g^{-1}y^\chi f_g$, with $f_g\in F_3'$.
In $\mathcal H$ this conjugator is central, so the action on
(E28.2) is exactly
\[
 (a,b,c)\longmapsto(\chi(g)a,\chi(g)b,\chi(g)^2c).
 \tag{E28.4}
\]
The unit tangent at $t=0$ lifts to the plus cubic by
$w=-x/z$, $t=w^3$ and $y/z=-(1-w^3)^{1/3}$.
Its toric leading coefficients are $(-1,-1)$.
Thus the base rational tangential section lifts to the cubic
and the inherited extension is $H\rtimes G_{\mathbb Q}$.
The pro-$3$ Kummer classes of both these signs, and of the
sign at the origin used in E.24, are zero. Changing between
these two tangential origins therefore fixes the toric
coordinates of every abelian class, so it leaves (E28.1)
and its entire fiber unchanged.

Reduce all three coordinates in (E28.3) modulo $27$.
This gives a Galois-stable finite quotient $H_{27}$ of order
$27^3$, of exponent $27$, whose commutator of lifts of
$(1,0)$ and $(0,1)$ is $(0,0,9)\ne1$.
At the local field $\mathbb Q_{109}$ the action (E28.4) on
$H_{27}$ is trivial: the $3$-adic cyclotomic character is
unramified there, with Frobenius $109\equiv1\pmod {27}$.
The integer $6$ has exact order $108$ in
$\mathbb F_{109}^{\times}$. Consequently its Kummer class
modulo $27$ is unramified with Frobenius value a generator of
$\mu_{27}$, while the Kummer class of $109$ has tame-inertia
value a generator. Choose the common generator of $\mu_{27}$
so that the Frobenius value for $6$ is $1$; a generator of
tame inertia can then be chosen with the value for $109$
also equal to $1$. We may choose Frobenius fixing the selected
roots of the uniformizer. The two projected values are thus
$(1,0)$ and $(0,1)$.

A section would therefore give a homomorphism
$G_{\mathbb Q_{109}}\to H_{27}$ with those two projections.
Wild inertia is pro-$109$ and maps trivially to this finite
$3$-group. In the tame quotient, choose Frobenius $\sigma$
and tame inertia $\tau$ using roots of the uniformizer.
Frobenius acts on their roots of unity by the $109$th power,
so
$\sigma\tau\sigma^{-1}=\tau^{109}$.
Since $H_{27}$ has exponent $27$, this relation forces the
two images to commute. Equation (E28.3) gives their commutator
as $(0,0,9)$, independently of their central coordinates,
a contradiction.
The program \texttt{e\_cubic\_class2.py} checks the order of
$6$ modulo $109$, the exponent and the full multiplication
formula, and all $27^2$ choices of the two central lifts.

Finally the map $N_3\to H_{27}$ kills $\gamma_3N_3$,
$N_3''$ and $\gamma_3 I_3$.
These are the actual factorizations through the quotients in
the statement. The explicit toric projection is $(a,b)$, so
every hypothetical lift of any point of the E.24 fiber would
give the forbidden homomorphism. E.24 proves the asserted
nonemptiness and size of the abelian fiber independently of
this obstruction. There is no reduction map from the rational
group $H_{\mathbb Q_3}$ to $H_{27}$ extending this reduction on
its integral lattice: every element of $H_{\mathbb Q_3}$ is a
$27$th power, whereas every $27$th power in $H_{27}$ is one.
Thus any such homomorphism is trivial.
The finite obstruction is used only for the indicated integral
pro-$3$ quotients. The result does not classify the remaining
cubic class-two sections whose toric pairs pass this local test.
\end{proof}
